%%%%%%%%%%%%%%%%%%%%%%%%%%%%%%%%%%%%%%%%%%%%%%%%%%%%%%%%%%%%%%%%%%%%%%%%%%%%%%%%
%% Plantilla de memoria en LaTeX para la EIF - Universidad Rey Juan Carlos
%%
%% Por Gregorio Robles <grex arroba gsyc.urjc.es>
%%     Grupo de Sistemas y Comunicaciones
%%     Escuela de Ingeniería de Fuenlabrada
%%     Universidad Rey Juan Carlos
%% (muchas ideas tomadas de Internet, colegas del GSyC, antiguos alumnos...
%%  etc. Muchas gracias a todos)
%%
%% La última versión de esta plantilla está siempre disponible en:
%%     https://github.com/gregoriorobles/plantilla-memoria
%%
%% Para obtener PDF, ejecuta en la shell:
%%   make
%% (las imágenes deben ir en PNG o JPG)

%%%%%%%%%%%%%%%%%%%%%%%%%%%%%%%%%%%%%%%%%%%%%%%%%%%%%%%%%%%%%%%%%%%%%%%%%%%%%%%%

\documentclass[a4paper, 12pt]{book}
%\usepackage[T1]{fontenc}

\usepackage[a4paper, left=2.5cm, right=2.5cm, top=3cm, bottom=3cm]{geometry}
\usepackage{times}
\usepackage[utf8]{inputenc}
\usepackage[spanish]{babel} % Comenta esta línea si tu memoria es en inglés
\usepackage{url}
%\usepackage[dvipdfm]{graphicx}
\usepackage{graphicx}
\usepackage{float}  %% H para posicionar figuras
\usepackage[nottoc, notlot, notlof, notindex]{tocbibind} %% Opciones de índice
\usepackage{latexsym}  %% Logo LaTeX

% Escribe el título y el nombre del autor / autora para que se use bien
% en otras partes de la plantilla
% Dependiendo de las partes de la plantilla, a veces aparecerán tal
% cual los escribas, a veces totalmente en mayúsculas, a veces de otras
% formas
\title{Caracterización de Commits de Git en Proyectos Software}
\author{Héctor Francisco Márquez García}

% Guarda el título, el autor y la fecha en variables
\makeatletter
\let\thetitle\@title
\let\theauthor\@author
\let\thedate\@date
\makeatother

\renewcommand{\baselinestretch}{1.5}  %% Interlineado

\begin{document}

\renewcommand{\refname}{Bibliografía}  %% Renombrando
\renewcommand{\appendixname}{Apéndice}


%%%%%%%%%%%%%%%%%%%%%%%%%%%%%%%%%%%%%%%%%%%%%%%%%%%%%%%%%%%%%%%%%%%%%%%%%%%%%%%%
% PORTADA

\begin{titlepage}
\begin{center}
\includegraphics[scale=0.6]{img/URJ_logo_Color_POS.png}

\vspace{1.75cm}

\LARGE
ESCUELA DE INGENIERÍA DE FUENLABRADA
\vspace{1cm}

\LARGE
GRADO EN INGENIERIA EN TECNOLOGIAS DE LA TELECOMUNICACIÓN 

\vspace{1cm}
\LARGE
\textbf{TRABAJO FIN DE GRADO}

\vspace{2cm}

\Large
\MakeUppercase{\thetitle}

\vspace{2cm}

\large
Autor : \theauthor \\
Tutor : Dr. Gregorio Robles Martínez\\
\vspace{1cm}

\large
Curso académico 2025/2026

\end{center}
\end{titlepage}

\newpage
\mbox{}
\thispagestyle{empty} % para que no se numere esta pagina



%%%%%%%%%%%%%%%%%%%%%%%%%%%%%%%%%%%%%%%%%%%%%%%%%%%%%%%%%%%%%%%%%%%%%%%%%%%%%%%%
%%%% Licencia
\clearpage
\pagenumbering{gobble}
\chapter*{}

\vspace{12cm}

%% Licencia de publicación en abierto elegida
%% Ver detalles en https://ofilibre.urjc.es/guias/tfg-abierto/

\begin{flushright}
\includegraphics[scale=0.6]{img/by-sa}
%\includegraphics[scale=0.6]{img/by}

%% Poner el año adecuado
\noindent©2026 \theauthor  \\
Algunos derechos reservados  \\
Este documento se distribuye bajo la licencia \\
``Atribución-CompartirIgual 4.0 Internacional'' de Creative Commons, \\
disponible en \\
\url{https://creativecommons.org/licenses/by-sa/4.0/deed.es}
\end{flushright}

%%%%%%%%%%%%%%%%%%%%%%%%%%%%%%%%%%%%%%%%%%%%%%%%%%%%%%%%%%%%%%%%%%%%%%%%%%%%%%%%
%%%% Dedicatoria

\chapter*{}
\pagenumbering{Roman} % para comenzar la numeracion de paginas en numeros romanos
\begin{flushright}
\textit{Dedicado a \\
mi familia / mi abuelo / mi abuela}
\end{flushright}

%%%%%%%%%%%%%%%%%%%%%%%%%%%%%%%%%%%%%%%%%%%%%%%%%%%%%%%%%%%%%%%%%%%%%%%%%%%%%%%%
%%%% Agradecimientos

\chapter*{Agradecimientos}
%\addcontentsline{toc}{chapter}{Agradecimientos} % si queremos que aparezca en el índice
\markboth{AGRADECIMIENTOS}{AGRADECIMIENTOS} % encabezado 

Aquí vienen los agradecimientos\ldots Aunque está bien acordarse de la pareja, no hay que olvidarse de dar las gracias a tu madre, que aunque a veces no lo parezca disfrutará tanto de tus logros como tú\ldots 
Además, la pareja quizás no sea para siempre, pero tu madre sí.

%%%%%%%%%%%%%%%%%%%%%%%%%%%%%%%%%%%%%%%%%%%%%%%%%%%%%%%%%%%%%%%%%%%%%%%%%%%%%%%%
%%%% Resumen

\chapter*{Resumen}
%\addcontentsline{toc}{chapter}{Resumen} % si queremos que aparezca en el índice
\markboth{RESUMEN}{RESUMEN} % encabezado

Aquí viene un resumen del proyecto.
Ha de constar de tres o cuatro párrafos, donde se presente de manera clara y concisa de qué va el proyecto. 
Han de quedar respondidas las siguientes preguntas:

\begin{itemize}
  \item ¿De qué va este proyecto? ¿Cuál es su objetivo principal?
  \item ¿Cómo se ha realizado? ¿Qué tecnologías están involucradas?
  \item ¿En qué contexto se ha realizado el proyecto? ¿Es un proyecto dentro de un marco general?
\end{itemize}

Lo mejor es escribir el resumen al final.

%%%%%%%%%%%%%%%%%%%%%%%%%%%%%%%%%%%%%%%%%%%%%%%%%%%%%%%%%%%%%%%%%%%%%%%%%%%%%%%%
%%%% Resumen en inglés

\chapter*{Summary}
%\addcontentsline{toc}{chapter}{Summary} % si queremos que aparezca en el índice
\markboth{SUMMARY}{SUMMARY} % encabezado

Here comes a translation of the ``Resumen'' into English. 
Please, double check it for correct grammar and spelling.
As it is the translation of the ``Resumen'', which is supposed to be written at the end, this as well should be filled out just before submitting.


%%%%%%%%%%%%%%%%%%%%%%%%%%%%%%%%%%%%%%%%%%%%%%%%%%%%%%%%%%%%%%%%%%%%%%%%%%%%%%%%
%%%%%%%%%%%%%%%%%%%%%%%%%%%%%%%%%%%%%%%%%%%%%%%%%%%%%%%%%%%%%%%%%%%%%%%%%%%%%%%%
% ÍNDICES %
%%%%%%%%%%%%%%%%%%%%%%%%%%%%%%%%%%%%%%%%%%%%%%%%%%%%%%%%%%%%%%%%%%%%%%%%%%%%%%%%

% Las buenas noticias es que los índices se generan automáticamente.
% Lo único que tienes que hacer es elegir cuáles quieren que se generen,
% y comentar/descomentar esa instrucción de LaTeX.

%%%% Índice de contenidos
\tableofcontents 
%%%% Índice de figuras
\cleardoublepage
%\addcontentsline{toc}{chapter}{Lista de figuras} % para que aparezca en el indice de contenidos
\listoffigures % indice de figuras
%%%% Índice de tablas
\cleardoublepage
%\addcontentsline{toc}{chapter}{Lista de tablas} % para que aparezca en el indice de contenidos
\listoftables % indice de tablas


%%%%%%%%%%%%%%%%%%%%%%%%%%%%%%%%%%%%%%%%%%%%%%%%%%%%%%%%%%%%%%%%%%%%%%%%%%%%%%%%
%%%%%%%%%%%%%%%%%%%%%%%%%%%%%%%%%%%%%%%%%%%%%%%%%%%%%%%%%%%%%%%%%%%%%%%%%%%%%%%%
% INTRODUCCIÓN %
%%%%%%%%%%%%%%%%%%%%%%%%%%%%%%%%%%%%%%%%%%%%%%%%%%%%%%%%%%%%%%%%%%%%%%%%%%%%%%%%

\cleardoublepage
\chapter{Introducción}
\label{sec:intro} % etiqueta para poder referenciar luego en el texto con ~\ref{sec:intro}
\pagenumbering{arabic} % para empezar la numeración de página con números

En las últimas décadas, el desarrollo de software colaborativo ha tenido un crecimiento significativo gracias al auge de los repositorios de código abierto y de plataformas de control de versiones distribuidas como Git. El software colaborativo es muy importante y utilizado en muchas industrias; todas estas compañías que desarrollan software necesitan tener un control sobre sus desarrollos colaborativos. Muchas de ellas utilizan Git como plataforma, sobre la cual tratará el estudio. 

Estos repositorios almacenan de forma detallada el historial evolutivo de los proyectos, abriendo la puerta a analizar cómo cambian los desarrollos a lo largo del tiempo mediante el estudio de los \emph{commits} realizados por los desarrolladores.

Siempre que un desarrollador envía su propuesta, capturando el estado del proyecto o segmento de este para el desarrollo del software, necesita realizar un \emph{commit} con los cambios. Cada \emph{commit} representa una unidad de cambio que puede reflejar decisiones de diseño, correcciones de errores, mejoras funcionales o tareas de mantenimiento. Por tanto, el análisis de los \emph{commits} puede proporcionar información relevante sobre los patrones de desarrollo, la evolución del software y las prácticas de ingeniería empleadas en cada proyecto. Es posible detectar y estudiar, de acuerdo a los históricos almacenados en plataformas como Git, si se ha introducido un error debido a la gran cantidad de cambios aplicados en un mismo \emph{commit}, o evaluar qué desarrolladores han tenido mayor impacto en el desarrollo.

Sin embargo, los repositorios de gran tamaño contienen una cantidad muy elevada de información, lo que dificulta su análisis manual. En este contexto, resulta necesario desarrollar herramientas automáticas capaces de extraer métricas, clasificar cambios y generar análisis empíricos que faciliten la comprensión de la evolución del software, para poder analizar cualquier posible falla o estudiar el comportamiento entre los desarrolladores.Por eso se decidió realizar este estudio y herramienta \emph{RepoAnalyzer}, la cual nos ayudara a categorizar y realizar un estudio que ayude a comprender el evolutivo de los desarrollos colaborativos.


\section{Contexto}
\label{sec:contexto}

\subsection{Repositorios de software libre}
\label{subsec:repositorios}

Los repositorios de software libre constituyen la base del desarrollo colaborativo moderno. Plataformas como GitHub o GitLab permiten almacenar el código fuente de los proyectos, gestionar versiones y coordinar el trabajo de múltiples desarrolladores distribuidos geográficamente. En un mundo donde se habitúa el trabajo y las colaboraciones desde diversos lugares del mundo, estas plataformas se han vuelto esenciales. Esto permite trabajar de forma continua en el desarrollo del software y aportar nuevas ideas o puntos de vista, aportando conocimiento de cualquier parte del planeta. Estas plataformas se manejan con unidades de cambio que se fusionan con el proyecto, estas unidades de cambio se denominan \emph{commit}.

Estos repositorios no solo contienen el código actual del proyecto, sino también todo su historial de cambios, lo que permite analizar su evolución a lo largo del tiempo. Esta característica resulta especialmente relevante en el ámbito de la ingeniería del software, ya que proporciona una fuente de datos rica para estudios empíricos sobre desarrollo colaborativo.

Además, muchos de estos proyectos son utilizados en entornos reales, desde librerías de programación hasta sistemas completos, lo que convierte su análisis en una oportunidad para comprender cómo se construye software a gran escala en contextos reales.

\subsection{Concepto de commit}
\label{subsec:commit}

Un \emph{commit} representa una unidad de cambio dentro de un repositorio de software. Cada \emph{commit} registra una serie de modificaciones realizadas sobre uno o varios ficheros, junto con información asociada como el autor, la fecha y un mensaje descriptivo de la finalidad de esa aportación.

Desde el punto de vista del desarrollo, los \emph{commits} pueden corresponder a diferentes tipos de actividades, como la corrección de errores, la incorporación de nuevas funcionalidades, tareas de mantenimiento o re-factorización del código. Esta diversidad de propósitos hace que el análisis de los mensajes de \emph{commit} sea especialmente útil para comprender la naturaleza de los cambios realizados, siempre y cuando sean organizados y no se mezclen diferentes finalidades en un mismo \emph{commit}.

Asimismo, el tamaño de un \emph{commit}, medido en número de ficheros modificados o líneas de código afectadas, puede proporcionar información relevante sobre la complejidad del cambio. En este contexto, resulta de interés identificar \emph{commits} que agrupan múltiples modificaciones independientes, conocidos como \emph{tangled commits}, ya que pueden dificultar el mantenimiento y la comprensión del código.

\subsection{Evolución del desarrollo colaborativo}
\label{subsec:evolucion}

El desarrollo de software colaborativo ha crecido de forma significativa en los últimos años, impulsado por la adopción de herramientas de control de versiones distribuidas y plataformas de colaboración en línea. Este crecimiento ha dado lugar a proyectos de gran escala, con miles de desarrolladores contribuyendo de manera simultánea.

Como consecuencia, los repositorios software actuales contienen grandes volúmenes de información histórica, lo que plantea nuevos retos en su análisis. Comprender cómo evolucionan estos proyectos, cómo se distribuye el trabajo entre los desarrolladores o qué patrones siguen los cambios realizados se ha convertido en un área de interés dentro de la ingeniería del software~\cite{hassan2008road}.

Diversos estudios han analizado la evolución de proyectos de software libre desde distintas perspectivas. Por ejemplo, trabajos como el de Mockus et al.~\cite{mockus2002two} estudian el desarrollo de grandes proyectos colaborativos, evidenciando diferencias significativas en la organización y dinámica de trabajo entre comunidades abiertas y entornos corporativos. Asimismo, investigaciones posteriores han profundizado en el análisis de historiales de \emph{commits} como fuente de información para comprender el comportamiento de los desarrolladores y la evolución del software.

En este contexto, el desarrollo de herramientas como \emph{RepoAnalyzer} permite abordar estos retos mediante el análisis automático de los historiales de \emph{commits}. Este tipo de herramientas facilita la obtención de métricas y la identificación de patrones que serían difíciles de detectar mediante análisis manual.



%%%%%%%%%%%%%%%%%%%%%%%%%%%%%%%%%%%%%%%%%%%%%%%%%%%%%%%%%%%%%%%%%%%%%%%%%%%%%%%%
%%%%%%%%%%%%%%%%%%%%%%%%%%%%%%%%%%%%%%%%%%%%%%%%%%%%%%%%%%%%%%%%%%%%%%%%%%%%%%%%
% OBJETIVOS %
%%%%%%%%%%%%%%%%%%%%%%%%%%%%%%%%%%%%%%%%%%%%%%%%%%%%%%%%%%%%%%%%%%%%%%%%%%%%%%%%

\cleardoublepage % empezamos en página impar
\chapter{Objetivos} % título del capítulo (se muestra)
\label{chap:objetivos} % identificador del capítulo (no se muestra, es para poder referenciarlo)

\section{Objetivo general} % título de sección (se muestra)
\label{sec:objetivo-general} % identificador de sección (no se muestra, es para poder referenciarla)

El objetivo principal de este trabajo fin de grado es diseñar e implementar una herramienta capaz de analizar repositorios de software basados en Git, con el fin de caracterizar los \emph{commits} y estudiar la evolución del desarrollo colaborativo en proyectos de código abierto.

Para ello, se pretende extraer y procesar información relevante del historial de cambios de los repositorios, permitiendo identificar patrones de desarrollo, analizar el comportamiento de los contribuidores y evaluar métricas relacionadas con la actividad del proyecto.

La herramienta desarrollada, denominada \emph{RepoAnalyzer}, tiene como finalidad automatizar este proceso de análisis, facilitando el estudio empírico de repositorios software reales y proporcionando una visión estructurada de su evolución.


\section{Objetivos específicos}
\label{sec:objetivos-especificos}

Para alcanzar el objetivo general planteado, se han definido los siguientes objetivos específicos:

\begin{itemize}

\item Analizar documentación técnica y estudios previos relacionados con la minería de repositorios software y su evolución en los desarrollos colaborativos , con el fin de identificar métricas, enfoques y metodologías relevantes aplicables al desarrollo de la herramienta.

\item Obtener información de repositorios Git a partir de su URL, permitiendo la descarga y el procesamiento de su historial de \emph{commits}, incluyendo autor, fecha, mensaje y métricas asociadas.

\item Diseñar un modelo de datos que represente de forma estructurada la información relevante de cada \emph{commit}, incluyendo autor, fecha, mensaje y métricas asociadas.

\item Implementar un sistema de clasificación semántica de \emph{commits} basado en palabras clave, que permita categorizar los cambios según su propósito.

\item Calcular métricas cuantitativas relacionadas con los \emph{commits}, como número de ficheros modificados, inserciones y eliminaciones.

\item Analizar la distribución de actividad entre los desarrolladores, identificando patrones de contribución dentro del repositorio.

\item Estudiar los intervalos temporales entre \emph{commits}, con el objetivo de caracterizar el ritmo de desarrollo del proyecto.

\item Detectar \emph{tangled commits}, es decir, cambios que agrupan múltiples modificaciones en un único \emph{commit}, y analizar su impacto.

\item Generar informes automáticos en distintos formatos (CSV, JSON y Excel) que faciliten la interpretación de los resultados obtenidos.

\end{itemize}


\section{Planificación temporal}
\label{sec:planificacion-temporal}
El desarrollo de este trabajo fin de grado se ha llevado a cabo a lo largo de varios meses, combinando su realización con la actividad laboral y el estudio académico. La mayor parte del trabajo se ha realizado fuera del horario laboral, incluyendo fines de semana y periodos vacacionales.

En una primera fase se realizó el análisis del problema y la revisión del estado del arte, con el objetivo de comprender las distintas aproximaciones existentes en el ámbito del análisis de repositorios software.

Posteriormente, se abordó el diseño e implementación de la herramienta \emph{RepoAnalyzer}, desarrollando de forma iterativa los distintos módulos encargados de la extracción, procesamiento y análisis de los datos.

Finalmente, se llevó a cabo la validación del sistema mediante la ejecución de la herramienta sobre distintos repositorios reales, así como la redacción de la memoria del trabajo.

Este enfoque progresivo ha permitido compaginar el desarrollo técnico del proyecto con la documentación del mismo, asegurando la coherencia entre los objetivos planteados y los resultados obtenidos.

%%%%%%%%%%%%%%%%%%%%%%%%%%%%%%%%%%%%%%%%%%%%%%%%%%%%%%%%%%%%%%%%%%%%%%%%%%%%%%%%
%%%%%%%%%%%%%%%%%%%%%%%%%%%%%%%%%%%%%%%%%%%%%%%%%%%%%%%%%%%%%%%%%%%%%%%%%%%%%%%%
% ESTADO DEL ARTE %
%%%%%%%%%%%%%%%%%%%%%%%%%%%%%%%%%%%%%%%%%%%%%%%%%%%%%%%%%%%%%%%%%%%%%%%%%%%%%%%%

\cleardoublepage
\chapter{Estado del arte}
\label{chap:estado}

En esta sección se mencionará el estado actual de las tecnologías sin las que no hubiera sido posible la ejecución de este trabajo.

\section{Mining Software Repositories}
\label{sec:msr}

El análisis de repositorios software, conocido como \emph{Mining Software Repositories} (MSR), es un área de investigación centrada en la extracción de conocimiento a partir de los datos históricos generados durante el desarrollo de software.

Los repositorios de control de versiones, sistemas de seguimiento de errores y plataformas de colaboración almacenan grandes volúmenes de información que pueden ser utilizados para estudiar el comportamiento entre los desarrolladores, la evolución del software y la calidad del código.

Hassan~\cite{hassan2008road} destaca que estos repositorios constituyen una fuente valiosa de información detallada, permitiendo analizar patrones de desarrollo y mejorar los procesos de ingeniería del software. Este enfoque ha dado lugar a numerosos estudios que utilizan métricas derivadas de los historiales de cambios para comprender mejor cómo se construyen y mantienen los desarrollos software.

\section{Clasificación semántica de commits}
\label{sec:clasificacion}

Uno de los aspectos más relevantes en el análisis de repositorios software es la clasificación de los \emph{commits} según su propósito. Los mensajes de \emph{commit} permiten identificar la intención del cambio realizado, como la corrección de errores, la implementación de nuevas funcionalidades o tareas de mantenimiento, en base de la clasificación que el desarrollador a establecido a cada \emph{commit}.

Diversos trabajos han propuesto enfoques para clasificar automáticamente los \emph{commits} utilizando técnicas basadas en palabras clave o métodos de aprendizaje automático. Estas clasificaciones permiten analizar la distribución de tipos de cambios dentro de un proyecto y estudiar la evolución de las actividades de desarrollo.

La clasificación semántica de los \emph{commits} resulta especialmente útil para comprender el comportamiento de los desarrolladores y la naturaleza de los cambios realizados a lo largo del tiempo. Con estos estudios se podrá dar una perspectiva de previsión de errores enfatizando la revisión del código cuando se observen ciertos patrones, además mejorar la carga de trabajo entre los desarrolladores.

\section{Tangled commits}
\label{sec:tangled}

Un problema identificado en el análisis de repositorios software es la existencia de \emph{tangled commits}, es decir, \emph{commits} que agrupan múltiples cambios lógicos en una única unidad.

Herzig et al.~\cite{herzig2013impact} analizan el impacto de este tipo de cambios en modelos de predicción de defectos, mostrando que los \emph{commits} pueden contener modificaciones heterogéneas que dificultan su interpretación y análisis automático, es decir, a la hora de clasificacion, un \emph{commit} puede venir asignado para una nueva funcionalidad, pero el \emph{commit} incluye además arreglos en otras zonas. Con este ultimo punto mencionado, se hace complicado la categorización de los \emph{commits} y su correcto estudio.

La presencia de \emph{tangled commits} introduce ruido en los estudios empíricos, ya que un único \emph{commit} puede incluir cambios de distinta naturaleza. Por ello, su identificación y análisis resulta relevante para mejorar la calidad de los estudios sobre evolución del software.

\section{Métricas y evolución del software}
\label{sec:metricas}

El estudio de la evolución del software se basa en el análisis de métricas derivadas de los historiales de cambios. Entre las métricas más utilizadas se encuentran el número de \emph{commits}, el tamaño de los cambios, el número de desarrolladores y la frecuencia de actividad.

Mockus et al.~\cite{mockus2002two} analizan proyectos de software libre como Apache y Mozilla, mostrando que la actividad de desarrollo suele concentrarse en un número reducido de contribuidores. Este tipo de análisis permite identificar patrones de colaboración y entender la dinámica de los desarrolladores en los proyectos.

Asimismo, el análisis temporal de los \emph{commits} permite estudiar el ritmo de desarrollo, identificar periodos de mayor actividad y comprender la evolución del proyecto a lo largo del tiempo.

\section{Limitaciones de los estudios existentes}
\label{sec:limitaciones}

A pesar de los avances en el análisis de repositorios software, muchos estudios se centran en aspectos específicos de forma aislada, como el tamaño de los \emph{commits} o la clasificación de cambios.

Esta fragmentación dificulta la obtención de una visión global de la evolución del software, ya que no se integran diferentes dimensiones de análisis en un mismo enfoque.

En este contexto, surge la necesidad de desarrollar herramientas que permitan combinar métricas cuantitativas, clasificación semántica y análisis temporal, proporcionando una visión más completa del comportamiento de los proyectos software.


%%%%%%%%%%%%%%%%%%%%%%%%%%%%%%%%%%%%%%%%%%%%%%%%%%%%%%%%%%%%%%%%%%%%%%%%%%%%%%%%
%%%%%%%%%%%%%%%%%%%%%%%%%%%%%%%%%%%%%%%%%%%%%%%%%%%%%%%%%%%%%%%%%%%%%%%%%%%%%%%%
% DISEÑO E IMPLEMENTACIÓN %
%%%%%%%%%%%%%%%%%%%%%%%%%%%%%%%%%%%%%%%%%%%%%%%%%%%%%%%%%%%%%%%%%%%%%%%%%%%%%%%%

\cleardoublepage
\chapter{Diseño e implementación}
\label{sec:diseno}

Aquí viene todo lo que has hecho tú (tecnológicamente). 
Puedes entrar hasta el detalle. 
Es la parte más importante de la memoria, porque describe lo que has hecho tú.
Eso sí, normalmente aconsejo no poner código, sino diagramas.



\section{Arquitectura general} 
\label{sec:arquitectura}

Si tu proyecto es un software, siempre es bueno poner la arquitectura (que es cómo se estructura tu programa a ``vista de pájaro'').

\begin{figure}
  \centering
  \includegraphics[width=9cm, keepaspectratio]{img/arquitectura.png}
  \caption{Estructura del parser básico}
  \label{fig:arquitectura}
\end{figure}


Por ejemplo, puedes verlo en la figura~\ref{fig:arquitectura}.
\LaTeX \ pone las figuras donde mejor cuadran. 
Y eso quiere decir que quizás no lo haga donde lo hemos puesto\ldots 
Eso no es malo.
A veces queda un poco raro, pero es la filosofía de \LaTeX: tú al contenido, que yo me encargo de la maquetación.


 
Recuerda que toda figura que añadas a tu memoria debe ser explicada.
Sí, aunque te parezca evidente lo que se ve en la figura~\ref{fig:arquitectura}, la figura en sí solamente es un apoyo a tu texto.
Así que explica lo que se ve en la figura, haciendo referencia a la misma tal y como ves aquí.
Por ejemplo: En la figura~\ref{fig:arquitectura} se puede ver que la estructura del \emph{parser} básico, que consta de seis componentes diferentes: los datos se obtienen de la red, y según el tipo de dato, se pasará a un \emph{parser} específico y bla, bla, bla\ldots

Si utilizas una base de datos, no te olvides de incluir también un diagrama de entidad-relación.


%%%%%%%%%%%%%%%%%%%%%%%%%%%%%%%%%%%%%%%%%%%%%%%%%%%%%%%%%%%%%%%%%%%%%%%%%%%%%%%%
%%%%%%%%%%%%%%%%%%%%%%%%%%%%%%%%%%%%%%%%%%%%%%%%%%%%%%%%%%%%%%%%%%%%%%%%%%%%%%%%
% EXPERIMENTOS Y VALIDACIÓN %
%%%%%%%%%%%%%%%%%%%%%%%%%%%%%%%%%%%%%%%%%%%%%%%%%%%%%%%%%%%%%%%%%%%%%%%%%%%%%%%%

\cleardoublepage
\chapter{Experimentos y validación}
\label{chap:experimentos}

Este capítulo se introdujo como requisito en 2019. 
Describe los experimentos y casos de test que tuviste que implementar para validar tus resultados. 
Incluye también los resultados de validación que permiten afirmar que tus resultados son correctos. 


%%%%%%%%%%%%%%%%%%%%%%%%%%%%%%%%%%%%%%%%%%%%%%%%%%%%%%%%%%%%%%%%%%%%%%%%%%%%%%%%
%%%%%%%%%%%%%%%%%%%%%%%%%%%%%%%%%%%%%%%%%%%%%%%%%%%%%%%%%%%%%%%%%%%%%%%%%%%%%%%%
% RESULTADOS %
%%%%%%%%%%%%%%%%%%%%%%%%%%%%%%%%%%%%%%%%%%%%%%%%%%%%%%%%%%%%%%%%%%%%%%%%%%%%%%%%

\cleardoublepage
\chapter{Resultados}
\label{chap:resultados}

En este capítulo se incluyen los resultados de tu trabajo fin de grado.

Si es una herramienta de análisis lo que has realizado, aquí puedes poner ejemplos de haberla utilizado para que se vea su utilidad.


%%%%%%%%%%%%%%%%%%%%%%%%%%%%%%%%%%%%%%%%%%%%%%%%%%%%%%%%%%%%%%%%%%%%%%%%%%%%%%%%
%%%%%%%%%%%%%%%%%%%%%%%%%%%%%%%%%%%%%%%%%%%%%%%%%%%%%%%%%%%%%%%%%%%%%%%%%%%%%%%%
% CONCLUSIONES %
%%%%%%%%%%%%%%%%%%%%%%%%%%%%%%%%%%%%%%%%%%%%%%%%%%%%%%%%%%%%%%%%%%%%%%%%%%%%%%%%

\cleardoublepage
\chapter{Conclusiones}
\label{chap:conclusiones}


\section{Consecución de objetivos}
\label{sec:consecucion-objetivos}

Esta sección es la sección espejo de las dos primeras del capítulo de objetivos, donde se planteaba el objetivo general y se elaboraban los específicos.

Es aquí donde hay que debatir qué se ha conseguido y qué no. 
Cuando algo no se ha conseguido, se ha de justificar, en términos de qué problemas se han encontrado y qué medidas se han tomado para mitigar esos problemas.

Y si has llegado hasta aquí, siempre es bueno pasarle el corrector ortográfico, que las erratas quedan fatal en la memoria final.
Para eso, en Linux tenemos aspell, que se ejecuta de la siguiente manera desde la línea de \emph{shell}:

\begin{verbatim}
  aspell --lang=es_ES -c memoria.tex
\end{verbatim}

\section{Aplicación de lo aprendido}
\label{sec:aplicacion}

Aquí viene lo que has aprendido durante el Grado/Máster y que has aplicado en el TFG/TFM.
Una buena idea es poner las asignaturas más relacionadas y comentar en un párrafo los conocimientos y habilidades puestos en práctica.

\begin{enumerate}
  \item a
  \item b
\end{enumerate}


\section{Lecciones aprendidas}
\label{sec:lecciones_aprendidas}

Aquí viene lo que has aprendido en el Trabajo Fin de Grado/Máster.

\begin{enumerate}
  \item Aquí viene uno.
  \item Aquí viene otro.
\end{enumerate}

\section{Declaración de Uso de Inteligencia Artificial}
\label{sec:uso_ia}

En este apartado se debería especificar el uso de inteligencia artificial, en particular, la IA generativa durante el TFG. No basta con decir ``usé IA". Debes especificar qué herramientas empleaste y para qué fases del trabajo y con qué propósito. Asimismo, debe quedar claro que, aunque la IA ayudó en el proceso, tú eres el autor único del contenido final y responsable de cualquier error u omisión.

A continuación se muestra un ejemplo de cómo se podría declarar el uso de IA en un trabajo:

\begin{itemize}
    \item \textbf{Herramientas utilizadas:} Se han empleado las herramientas \textit{[p.ej., ChatGPT-4o de OpenAI, Claude 3 de Anthropic y DeepL]}.
    
    \item \textbf{Alcance de la aplicación:} El uso de dichas herramientas se ha limitado exclusivamente a las siguientes tareas auxiliares:

    \begin{enumerate}
        \item Corrección gramatical, ortográfica y mejora de la fluidez estilística del texto original en español e inglés.
        \item Generación de ideas preliminares para la estructuración del índice y capítulos.
        \item Traducción de fragmentos de bibliografía técnica del inglés al español.
        \item Optimización y depuración de fragmentos de código en lenguaje Python.
    \end{enumerate}

    \item \textbf{Responsabilidad y Autoría:} El autor manifiesta que la IA ha actuado únicamente como un asistente de apoyo. Toda la argumentación, el análisis crítico, el diseño experimental y las conclusiones finales son de elaboración propia y original. 
    
    \item \textbf{Verificación de contenido:} Se garantiza que toda la información, datos y referencias bibliográficas generadas o sugeridas por las herramientas de IA han sido verificadas manualmente mediante fuentes académicas primarias para evitar errores o falsedades (alucinaciones).
\end{itemize}


\section{Trabajos futuros}
\label{sec:trabajos_futuros}

Ningún proyecto ni software se termina, así que aquí vienen ideas y funcionalidades que estaría bien tener implementadas en el futuro.

Es un apartado que sirve para dar ideas de cara a futuros TFGs/TFMs.


%%%%%%%%%%%%%%%%%%%%%%%%%%%%%%%%%%%%%%%%%%%%%%%%%%%%%%%%%%%%%%%%%%%%%%%%%%%%%%%%
%%%%%%%%%%%%%%%%%%%%%%%%%%%%%%%%%%%%%%%%%%%%%%%%%%%%%%%%%%%%%%%%%%%%%%%%%%%%%%%%
% APÉNDICE(S) %
%%%%%%%%%%%%%%%%%%%%%%%%%%%%%%%%%%%%%%%%%%%%%%%%%%%%%%%%%%%%%%%%%%%%%%%%%%%%%%%%

\cleardoublepage
\appendix
\chapter{Manual de usuario}
\label{app:manual}

Esto es un apéndice.
Si has creado una aplicación, siempre viene bien tener un manual de usuario.
Pues ponlo aquí.

%%%%%%%%%%%%%%%%%%%%%%%%%%%%%%%%%%%%%%%%%%%%%%%%%%%%%%%%%%%%%%%%%%%%%%%%%%%%%%%%
%%%%%%%%%%%%%%%%%%%%%%%%%%%%%%%%%%%%%%%%%%%%%%%%%%%%%%%%%%%%%%%%%%%%%%%%%%%%%%%%
% BIBLIOGRAFIA %
%%%%%%%%%%%%%%%%%%%%%%%%%%%%%%%%%%%%%%%%%%%%%%%%%%%%%%%%%%%%%%%%%%%%%%%%%%%%%%%%

\cleardoublepage

% Las siguientes dos instrucciones es todo lo que necesitas
% para incluir las citas en la memoria
\bibliographystyle{abbrv}
\bibliography{memoria}  % memoria.bib es el nombre del fichero que contiene
% las referencias bibliográficas. Abre ese fichero y mira el formato que tiene,
% que se conoce como BibTeX. Hay muchos sitios que exportan referencias en
% formato BibTeX. Prueba a buscar en http://scholar.google.com por referencias
% y verás que lo puedes hacer de manera sencilla.
% Más información: 
% http://texblog.org/2014/04/22/using-google-scholar-to-download-bibtex-citations/

\end{document}
